
\documentclass{article}
\usepackage{colortbl}
\usepackage{makecell}
\usepackage{multirow}
\usepackage{supertabular}

\begin{document}

\newcounter{utterance}

\centering \large Interaction Transcript for game `cladder', experiment `full\_v1.5\_default', episode 8201 with Qwen3.5{-}9B.
\vspace{24pt}

{ \footnotesize  \setcounter{utterance}{1}
\setlength{\tabcolsep}{0pt}
\begin{supertabular}{c@{$\;$}|p{.15\linewidth}@{}p{.15\linewidth}p{.15\linewidth}p{.15\linewidth}p{.15\linewidth}p{.15\linewidth}}
    \# & \multicolumn{2}{c}{Player} && \multicolumn{2}{c}{Game Master} \\
    \hline

    \theutterance \stepcounter{utterance}  
    & & & \multicolumn{4}{p{0.6\linewidth}}{
        \cellcolor[rgb]{0.9,0.9,0.9}{
            \makecell[{{p{\linewidth}}}]{
                \texttt{\tiny{[P1$\langle$GM]}}
                \texttt{You are an expert in causal inference. The following question is not a typical commonsense query, but rather a meticulously designed question created by a professor specializing in causal inference, intended to assess the students' mastery of the course content.} \\
\texttt{Imagine a self{-}contained, hypothetical world with only the following conditions, and without any unmentioned factors or causal relationships: Tijv has a direct effect on tijw and xevo. Xevo has a direct effect on gyzp. Tijw has a direct effect on gyzp. For those who are not xevo and are not tijw, the probability of gyzp is 76\%. For those who are not xevo and are tijw, the probability of gyzp is 76\%. For those who are xevo and are not tijw, the probability of gyzp is 12\%. For those who are xevo and are tijw, the probability of gyzp is 14\%. For those who are not xevo, the probability of tijw is 28\%. For those who are xevo, the probability of tijw is 14\%. For those who are xevo, would it be more likely to see gyzp if the individual was not xevo?} \\
\texttt{Start your answer with 'Yes' or 'No', followed by additional reasoning or evidence to support your explanation.} \\
            }
        }
    }
     \\ \\

    \theutterance \stepcounter{utterance}  
    & \multicolumn{4}{p{0.6\linewidth}}{
        \cellcolor[rgb]{0.561,0.737,0.561}{
            \makecell[{{p{\linewidth}}}]{
                \texttt{\tiny{[P1$\rangle$GM]}}
                \texttt{No} \\
\\ 
\texttt{**Reasoning:**} \\
\\ 
\texttt{To answer this question, we must compare the probability of observing \$gyzp\$ under two specific counterfactual scenarios for an individual who is currently \$xevo\$:} \\
\texttt{1.  **Actual State:** The individual is \$xevo\$.} \\
\texttt{2.  **Counterfactual State:** The individual is \$xevo\$ but is intervened upon to be "not \$xevo\$."} \\
\\ 
\texttt{Let \$P(gyzp | do(X))\$ denote the probability of \$gyzp\$ given the intervention on \$xevo\$. The question asks if \$P(gyzp | do(\textbackslash{}neg xevo)) > P(gyzp | do(xevo))\$.} \\
\\ 
\texttt{**Step 1: Determine \$P(gyzp | do(xevo))\$**} \\
\texttt{In the actual world where the individual is \$xevo\$, the probability of \$gyzp\$ is given directly by the problem statement:} \\
\texttt{"For those who are xevo and are tijw, the probability of gyzp is 14\%."} \\
\texttt{"For those who are xevo and are not tijw, the probability of gyzp is 12\%."} \\
\\ 
\texttt{We must calculate the total probability for the \$xevo\$ group by weighting these outcomes by the probability of being \$tijw\$ given \$xevo\$.} \\
\texttt{*   \$P(tijw | xevo) = 14\textbackslash{}\% = 0.14\$} \\
\texttt{*   \$P(\textbackslash{}neg tijw | xevo) = 1 {-} 0.14 = 0.86\$} \\
\\ 
\texttt{\$\$P(gyzp | xevo) = P(gyzp | xevo, tijw)P(tijw | xevo) + P(gyzp | xevo, \textbackslash{}neg tijw)P(\textbackslash{}neg tijw | xevo)\$\$} \\
\texttt{\$\$P(gyzp | xevo) = (0.14 \textbackslash{}times 0.14) + (0.12 \textbackslash{}times 0.86)\$\$} \\
\texttt{\$\$P(gyzp | xevo) = 0.0196 + 0.1032 = 0.1228 \textbackslash{}text\{ (or 12.28\textbackslash{}\%)\}\$\$} \\
\\ 
\texttt{**Step 2: Determine \$P(gyzp | do(\textbackslash{}neg xevo))\$**} \\
\texttt{If we intervene to set \$xevo = \textbackslash{}neg xevo\$, the causal mechanism changes. We must look at the probabilities defined for the "not \$xevo\$" group.} \\
\texttt{The problem states:} \\
\texttt{"For those who are not xevo and are not tijw, the probability of gyzp is 76\%."} \\
\texttt{"For those who are not xevo and are tijw, the probability of gyzp is 76\%."} \\
\\ 
\texttt{Crucially, in this hypothetical world, the probability of \$gyzp\$ is **76\%**} \\
            }
        }
    }
    & & \\ \\

    \theutterance \stepcounter{utterance}  
    & & & \multicolumn{2}{p{0.3\linewidth}}{
        \cellcolor[rgb]{0.9,0.9,0.9}{
            \makecell[{{p{\linewidth}}}]{
                \texttt{\tiny{[GM$|$GM]}}
                \texttt{No} \\
            }
        }
    }
    & & \\ \\

    \theutterance \stepcounter{utterance}  
    & & & \multicolumn{2}{p{0.3\linewidth}}{
        \cellcolor[rgb]{0.9,0.9,0.9}{
            \makecell[{{p{\linewidth}}}]{
                \texttt{\tiny{[GM$|$GM]}}
                \texttt{yes} \\
            }
        }
    }
    & & \\ \\

    \theutterance \stepcounter{utterance}  
    & & & \multicolumn{2}{p{0.3\linewidth}}{
        \cellcolor[rgb]{0.9,0.9,0.9}{
            \makecell[{{p{\linewidth}}}]{
                \texttt{\tiny{[GM$|$GM]}}
                \texttt{game\_result = LOSE} \\
            }
        }
    }
    & & \\ \\

\end{supertabular}
}

\end{document}
